\documentclass[12pt]{article}
\usepackage[utf8]{inputenc}
\usepackage[T1]{fontenc}
\usepackage{hyperref}
\title{Hyolmo Language Resources for K-12 Curriculum Development}
\author{Compiled by Claude Sonnet 5.5 (claude-sonnet-5-5) for Andr\'as Kornai}
\date{}
\begin{document}
\maketitle

\section{Introduction}
Hyolmo (Yolmo, Helambu Sherpa; ISO 639-3 scp) is a Tibetic language of Nepal. The 2011 census counted 10,752 Hyolmo people, of whom 10,176 gave Hyolmo as their mother tongue. Speakers live mainly in the Helambu and Melamchi valleys. This dossier lists only resources that appeared in web searches, and little school material was found.

\section{Primary Educational Resources}
No primers, textbooks or school programmes for Hyolmo were found. The Endangered Languages Project lists some community materials (2025), including a teaching video by a Yolmo lopen, but their details were not confirmed.
\begin{itemize}
\item Endangered Languages Project, Hyolmo resources page. \url{https://www.endangeredlanguages.com/node/39970/resources}
\end{itemize}

\section{Linguistic and Reference Materials}
\begin{itemize}
\item Anna Maria Hari and Chhegu Lama, \emph{Yolmo-Nepali-English Dictionary} (Tribhuvan University, Kathmandu, 2004, about 848 pages); Glottolog record. \url{https://glottolog.org/resource/reference/id/151906}
\item Second Glottolog record for Hari and Lama 2004. \url{https://glottolog.org/resource/reference/id/678332}
\item Wikibooks, Research on Hyolmo (grammatical description). \url{https://en.wikibooks.org/wiki/Research_on_Hyolmo}
\item Yolmo language overview. \url{https://en.wikipedia.org/wiki/Yolmo_language}
\item Hyolmo numeral system (Max Planck Institute). \url{https://lingweb.eva.mpg.de/channumerals/Hyolmo.htm}
\item Omniglot page on Yolmo, noting dictionaries in Devanagari and Latin script. \url{https://www.omniglot.com/writing/yolmo.htm}
\item SOAS repository item on Yolmo. \url{https://soas-repository.worktribe.com/OutputFile/344080}
\end{itemize}

\section{Religious and Community Resources}
\begin{itemize}
\item ``Who is Hyolmo?'' from Indigenous Voice. \url{https://english.indigenousvoice.com/hyolmo}
\item Yolmo people overview. \url{https://en.wikipedia.org/wiki/Yolmo_people}
\end{itemize}
No scripture translations were found.

\section{Digital and Audio Resources}
Videos listed on the Endangered Languages Project (see above) include ``Growing Hope in Yolmo''. No corpus or software was confirmed.

\section{Curriculum Development Context}
The language is spoken by almost all community members, but monolingual speakers are declining and migration to Kathmandu and abroad favours Nepali and English. Hyolmo has no established school orthography; both Devanagari and Latin systems are used in dictionaries. A curriculum would start from the 2004 dictionary and the descriptive grammar work, with new primers and readers written with the community.
\end{document}